\itemhead{3.1 思考 · 未编号题组（教材p.65）}{XJ-B1-C03-TH-P065-Y1561}
\label{XJ-B1-C03-TH-P065-Y1561}
{\footnotesize 湘教版教材p.65 / PDF第72页 / 3.1 思考 / 未编号题组}
\field{原题}{我们有了集合的概念，就可以将上述例子中两个变量的取值范围看成集合，将函数看成两个集合之间的对应关系。那么，上述例子中的函数有哪些共同特征？
\par\textit{教材本页所引用的例子：}$s=x+3$；梯形两底 $a,b$ 为常量时 $S=\dfrac{(a+b)h}2$；圆周长 $C=2\pi r$；自由落体路程 $H=\dfrac12gt^2\ (t\in[0,T])$。}
\field{参考答案}{每个关系都有明确的自变量取值集合，对其中每一个值，计算规则都给出唯一的实数结果；共同本质是非空数集之间的确定对应，而不只是一条公式。}
{\small \field{关联知识}{F01 函数本质、构成要素与相同函数（教材帮p.77、78 / PDF86、87）}
\field{方法与易错点}{逐一核对每个允许输入是否有且仅有一个合乎目标集合的输出。 易错：把值域直接当作目标集合；只看解析式外观；认为常函数不是函数。}
\field{答案依据}{答案册未提供本栏目的独立对应答案；依据原题、教材概念补写参考回应。开放性总结不设唯一数值答案。}}
